\documentclass[pdf, unicode, 12pt, a4paper, oneside, fleqn]{article}

% === Содержимое log-style ===
\setcounter{secnumdepth}{0}
\usepackage{float}

\usepackage[utf8]{inputenc}
\usepackage[T2B]{fontenc}
\usepackage[english,russian]{babel}

\frenchspacing

\usepackage{amsmath}
\usepackage{amssymb}
\usepackage{hyperref}
\usepackage{longtable}
\usepackage[table]{xcolor}
\usepackage{array}
\usepackage{color}
\usepackage{xcolor}

\newcommand{\MYhref}[3][blue]{\href{#2}{\color{#1}{#3}}}

\usepackage{listings}
\usepackage{alltt}
\usepackage{csquotes}
\DeclareQuoteStyle{russian}
	{\guillemotleft}{\guillemotright}[0.025em]
	{\quotedblbase}{\textquotedblleft}
\ExecuteQuoteOptions{style=russian}

\usepackage{graphicx}

\usepackage{minted}
\definecolor{bg}{rgb}{0.85,0.85,0.85}
\setminted{
    breaklines=true,
    breakanywhere=true,
    frame=none,
    bgcolor=bg,
    fontsize=\small
}

\lstset{
	tabsize=2,
	breaklines,
	columns=fullflexible,
	flexiblecolumns,
	numbers=left,
	numberstyle={\footnotesize},
	extendedchars,
	inputencoding=utf8
}

\def\@xobeysp{ }
\def\verbatim@processline{\hspace{1.2cm}\raggedright\the\verbatim@line\par}

\oddsidemargin=-0.4mm
\textwidth=160mm
\topmargin=4.6mm
\textheight=210mm

\parindent=0pt
\parskip=3pt

\definecolor{lightgray}{gray}{0.9}

\renewcommand{\thesubsection}{\arabic{subsection}}

\lstdefinestyle{customc}{
  belowcaptionskip=1\baselineskip,
  breaklines=true,
  frame=L,
  xleftmargin=\parindent,
  language=C,
  showstringspaces=false,
  basicstyle=\footnotesize\ttfamily,
  keywordstyle=\bfseries\color{green!40!black},
  commentstyle=\itshape\color{gray},
  identifierstyle=\color{black},
  stringstyle=\color{blue},
}

\lstdefinestyle{customasm}{
  belowcaptionskip=1\baselineskip,
  frame=L,
  xleftmargin=\parindent,
  language=[x86masm]Assembler,
  basicstyle=\footnotesize\ttfamily,
  commentstyle=\itshape\color{purple!40!black},
}

\lstset{escapechar=@,style=customc}

\newcommand{\CWPHeader}[1]{\addtocounter{section}{-1}\section{#1}}
\newcommand{\CWHeader}[1]{\section*{#1}}
\newcommand{\CWProblem}[1]{\par\textbf{Задача: }#1}
% === Конец log-style ===

\begin{document}

\begin{titlepage}
\begin{center}
\bfseries

{\Large Московский авиационный институт\\ (национальный исследовательский университет)}

\vspace{48pt}

{\large Факультет информационных технологий и прикладной математики}

\vspace{36pt}

{\large Кафедра вычислительной математики и~программирования}

\vspace{48pt}

Лабораторные работы по курсу \enquote{Информационный поиск}

\end{center}

\vspace{72pt}

\begin{flushright}
\begin{tabular}{rl}
Студент: & К.\,А. Былькова \\
Преподаватель: & А.\,А. Кухтичев \\
Группа: & М8О-408Б-22 \\
Дата: & \\
Оценка: & \\
Подпись: & \\
\end{tabular}
\end{flushright}

\vfill

\begin{center}
\bfseries
Москва, \the\year
\end{center}
\end{titlepage}

\pagebreak

\tableofcontents

\

\textbf{Ссылка на GitHub: \href{https://github.com/kr1st1na0/InfSearch_MAI}{https://github.com/kr1st1na0/InfSearch\_MAI}}

\pagebreak

\section{Лабораторная работа \textnumero 1 \enquote{Добыча корпуса документов}}

Необходимо подготовить корпус документов, который будет использован при выполнении остальных лабораторных работ:
\begin{itemize}
    \item Скачать его к себе на компьютер. В отчёте нужно указать источник данных.
    \item Ознакомиться с ним, изучить его характеристики. Из чего состоит текст? Есть ли
дополнительная мета-информация? Если разметка текста, какая она?
    \item Разбить на документы.
    \item Выделить текст.
    \item Найти существующие поисковики, которые уже можно использовать для поиска по
выбранному набору документов (встроенный поиск Википедии, поиск Google с использованием ограничений на URL или на сайт). Если такого поиска найти невозможно, то использовать корпус для выполнения лабораторных работ нельзя!
    \item Привести несколько примеров запросов к существующим поисковикам, указать недостатки в полученной поисковой выдаче.
\end{itemize}

В результатах работы должна быть указаны статистическая информация о корпусе:
\begin{itemize}
    \item Размер \enquote{сырых} данных.
    \item Количество документов.
    \item Размер текста, выделенного из \enquote{сырых} данных.
    \item Средний размер документа, средний объём текста в документе.
\end{itemize}

\subsection*{Описание}

\textbf{- Тема}

В качестве темы документов был выбран «Научпоп» (научно
популярные статьи по медицине, технологиям, науке, гаджетам и прочее). Выбор обусловлен тем, что в современном мире научно-популярные статьи очень распространены и помогают людям получать ценную информацию в доступном и простом формате. Немаловажным критерием при выборе также выступало наличие большого количества статей, связанных с данной тематикой. 

\textbf{- Источники данных}

Для создания корпуса документов были выбраны два источника на 
русском языке: \href{https://biomolecula.ru}{Biomolecula} и \href{https://www.techinsider.ru}{Techinsider}. 

\subsection*{Ход работы}

\textbf{- Проверка источников}

Для того, чтобы убедиться, что выбранные источники подходят для 
корпуса документов для выполнения лабораторных работ, была 
произведена проверка в поиске Яндекс (с использованием ограничений на 
URL или на сайт). 

Для проверки в поисковую строку были введены запросы вида 
site:<источник> <интересующий запрос>. Оба источника успешно прошли 
проверку: 

\begin{figure}[H]
    \centering
    \includegraphics[width=0.9\linewidth]{yandex_bio.png}
    \caption{Проверка источника Biomolecula}
    \label{fig:placeholder}
\end{figure}

\begin{figure}[H]
    \centering
    \includegraphics[width=0.9\linewidth]{yandex_techin.png}
    \caption{Проверка источника Techinsider}
    \label{fig:placeholder}
\end{figure}

\textbf{- Характеристики документов}

\begin{itemize}
    \item Biomolecula:
    \begin{itemize}
        \item Мета-информация: находится в <head>: <title>… </title>– название статьи, <meta content="…" name="author"> – автор статьи, <meta name=”description” – краткое описание статьи. В <body> можно найти дату публикации статьи: <p class="article\_date">20 ноября 2015</p>.
        \item Текст: в <body> в <div class="mainbar mainbar\_\_insert"> расположен текст статьи <p><em>. Заголовок раздела – <h2>…</h2>. Абзац – <p>…</p>.
    \end{itemize}
    \item Techinsider:
        \begin{itemize}
            \item Метаданные статьи (заголовок, автор и дата публикации) находятся в структурированных данных в формате JSON-LD, размещённых на странице внутри тегов <script type="application/ld+json">: заголовок в поле "headline"; автор в поле "author", которое может содержать как строку, так и объект с ключом "name"; дата публикации в поля "datePublished".
            \item Текст статьи также размещён не в HTML-абзацах, а внутри блока <script type="application/ld+json">. В этом JSON-блоке в поле articleBody содержится готовый к использованию чистый текст статьи, без HTML-разметки, рекламы или навигационных элементов.
        \end{itemize}
\end{itemize}

\textbf{- Пример отрывка из выделенного текста}

На протяжении многих веков лыжи представляли собой длинные (больше 2 м) плоские полозья с загибом спереди, вырезанные из сплошного куска дерева, обычно гикори. Лыжа шириной 10 см и толщиной 3 см весила больше двух килограммов. Изгиб лыжи вверх в середине (весовой прогиб) был сделан для равномерного распределения давления от лыжника по всей длине (а не только в зоне под ботинком). После нескольких спусков деревянную лыжу нужно было смазывать воском, чтобы к ней не прилипал снег. От катания по жесткому снегу и льду ребра быстро стачивались и расщеплялись. Хотя лыжи старались защищать от влаги, вода все равно попадала внутрь, и со временем лыжа скручивалась и теряла форму. Первые слоеные деревянные лыжи появились еще в 1890-х годах. Они были легче, чем лыжи из массива дерева, и стали особенно популярны в конце 1930-х годов, когда появились водостойкие клеи, с помощью которых соединяли слои шпона. К 1951 году 90\% производимых лыж были слоеными, но и они были далеки от идеала…

\textbf{- Информация о корпусе документов}

В результате был собран корпус из 106616 документов, включающий данные из двух источников: biomolecula.ru и techinsider.ru.

Статистика корпуса:

\begin{center}
\begin{tabular}{ |c|c| }
\hline
 Общее количество документов & 106616 документов \\ 
 \hline
 Общий размер «сырых» данных & 27 ГБ \\
 \hline
 Средний размер «сырых» данных & 267 КБ \\
 \hline
 Общий размер текста & 502 МБ \\
 \hline
 Средний размер текста & 4.8 КБ \\
 \hline
\end{tabular}
\end{center}

\subsection*{Выводы}

В результате выполнения данной лабораторной был собран уникальный корпус документов из двух источников по теме «Научпоп», пригодный для дальнейшей работы. Данный корпус достаточно большой, был проиндексирован поисковиком Яндекс. Была разобрана структура документов из каждого источника, найдена мета-информация.

\pagebreak

\section{Лабораторная работа \textnumero 2 \enquote{Поисковой робот}}

Необходимо написать парсер на любом языке программирования.
\begin{itemize}
    \item Написать поисковый робот — компоненты обкачки документов, 
используя любой язык программирования.
    \item Единственным аргументом поисковому роботу подаётся путь до yaml конфига, содержащий:
    \begin{itemize}
        \item Данные для базы данные в секции db;
        \item Данные для робота в секции logic: задержка между обкачкой страницы;
        \item Любые другие данные, необходимые для реализовации логики поискового робота.
    \end{itemize}
    \item Сохранять в базе данных (например, MongoDB) документы со следующими полями: 
    \begin{itemize}
        \item url, нормализованный;
        \item «сырой» html-текст документа;
        \item название источника; 
        \item дата обкачки документа в формате Unix time stamp.
    \end{itemize}
    \item Поисковый робот можно остановить в любой момент и при повторном 
запуске робот должен начать с того документа, с которого он 
остановился.
    \item Периодически он должен уметь переобкачивать документы, которые 
уже есть в базе, но только в том случае, если они изменились. 
\end{itemize}

\subsection*{Описание}

\textbf{- Язык программирования}

Поисковой робот реализован на языке Python.

\textbf{- База данных}

Данные хранятся в базе данных PostgreSQL. Таблица documents 
содержит следующие поля: 

\begin{itemize}
    \item url — исходный URL документа
    \item normalized\_url — нормализованный URL
    \item html — HTML-контент страницы
    \item clean\_text — извлечённый и очищенный текст статьи
    \item title, author, publish\_date — метаданные статьи
    \item source — название источника («Биомолекула» или «Техинсайдер»)
    \item fetch\_timestamp — время первоначального сохранения
    \item last\_modified\_header, etag\_header — HTTP-заголовки для определения изменений
    \item last\_fetch\_attempt — время последней попытки обновления документа
\end{itemize}

\textbf{- Config.yaml}

\begin{itemize}
    \item db:
    \begin{itemize}
        \item host: "localhost"
        \item port: 5432
        \item database: "docs\_corpus"
        \item user: "postgres"
    \end{itemize}
    \item logic:
    \begin{itemize}
        \item delay\_seconds: 1
        \item user\_agent: "MAI-InfoSearch (university homework; contact: bylkovaka@list.ru)"
        \item respect\_robots\_txt: true
        \item max\_documents: 120000
        \item sources:
        \begin{itemize}
            \item name: "Биомолекула"
            \item start\_urls: "https://biomolecula.ru/sitemap.xml"
            \item name: "Техинсайдер"
            \item start\_urls: "https://www.techinsider.ru/sitemap.xml"
        \end{itemize}
    \end{itemize}

\end{itemize}

\textbf{- Обход сайтов}

Обход осуществляется на основе карт сайтов (sitemap.xml). Для 
каждого источника («Биомолекула» и «Техинсайдер») загружаются URL из 
корневого sitemap-файла, включая вложенные sitemap. Рекурсивно 
извлекаются все ссылки на статьи и добавляются их в очередь обработки. 

\textbf{- Фильтрация URL}

Для каждого сайта заданы правила фильтрации, исключающие 
служебные страницы (например, категории, теги, профили пользователей, 
страницы поиска и т.п.). На «Техинсайдере» дополнительно требуются 
цифры в пути URL как признак статьи. На «Биомолекуле» исключаются все 
URL, начинающиеся с /tag/, /category/, /authors/ и аналогичных. 

\textbf{- Нормализация URL}

Все URL приводятся к единому каноническому виду: схема (https) и 
сетевое имя (netloc) — в нижнем регистре, путь — нормализован (удалены 
параметры запроса и фрагменты, добавлен корневой слэш при 
необходимости). Это позволяет корректно определять дубликаты, 
возникающие из-за редиректов или различий в регистре. 

\textbf{- Сохранение состояния}

Во время обхода периодически сохраняется текущее состояние в 
файл crawler\_state.pkl с использованием модуля pickle. В состояние входят:
\begin{itemize}
    \item очередь оставшихся URL-адресов для обработки
    \item множество уже добавленных в очередь URL (для предотвращения 
дублирования)
    \item количество уже сохранённых документов
\end{itemize}
Это позволяет прервать работу и возобновить её позже без потери прогресса. 

\textbf{- Обновление}

После первоначального сбора документов запускается фаза 
обновления. Робот выбирает из базы документы, последняя попытка 
обновления которых была более недели назад (last\_fetch\_attempt < now - 7 дней). Для каждого такого документа выполняется условный HTTP-запрос с заголовками If-Modified-Since и/или If-None-Match. 
\begin{itemize}
    \item если сервер возвращает статус 304 Not Modified, документ помечается как проверенный, содержимое не меняется
    \item если содержимое изменилось (статус 200), извлекаются обновлённые текст и метаданные, вычисляется хеш документа. При изменении хеша документ обновляется в базе
    \item поддерживается корректная обработка редиректов: если URL изменился, проверяется, не существует ли уже документ по новому normalized\_url, и при необходимости обновление отменяется, чтобы избежать дублирования
\end{itemize}

\subsection*{Ход работы}

\textbf{- Обход документов}

Обход начинается с загрузки карты сайта (sitemap.xml) каждого из 
двух целевых ресурсов — «Биомолекула» и «Техинсайдер». Робот 
рекурсивно обрабатывает вложенные sitemap-файлы, если они содержат 
другие sitemap-ссылки (например, biography1-utf8.xml у Techinsider), и 
извлекает из них все <loc>-элементы с URL статей. 

Каждый полученный URL проходит валидацию с помощью функции 
qualifies\_as\_article\_page, которая отсеивает служебные страницы по пути (/tag/, /category/, /authors/ и т.д.) и, в случае Techinsider, требует наличия цифр в URL как признака статьи.

Далее URL нормализуется и добавляется в очередь обработки, если 
он ещё не сохранён в базе и не находится в кэше seen\_in\_queue. 

\textbf{- Извлечение мета-информации}

Метаданные — заголовок, автор, дата публикации — извлекаются 
специфически для каждого источника с помощью функции 
parse\_article\_metadata. 
\begin{itemize}
    \item для Techinsider используется структурированный JSON-LD, встроенный в <script type="application/ld+json">
    \item для Биомолекулы данные берутся напрямую из HTML: заголовок из 
    <h1>, автор — из ссылки с классом author-link, дата — либо из тега <time datetime="...">, либо из текста вида «31 октября 2025», который парсится с помощью регулярного выражения и маппинга месяцев
\end{itemize}
Если нужные элементы не найдены, поля остаются пустыми, но сам 
документ всё равно может быть сохранён, если текст извлечён.

\textbf{- Извлечение текста}

Текст статьи очищается от навигации, рекламы, сносок и служебных 
блоков. Процесс сильно отличается по источникам: 
\begin{itemize}
    \item для Techinsider текст берётся напрямую из поля articleBody внутри JSON-LD; 
    \item для Биомолекулы ищется контейнер .mainbar\_\_insert, из которого удаляются нерелевантные блоки по CSS-селекторам. Затем извлекаются абзацы <p>, фильтруются по длине и содержанию (отсеиваются подписи рисунков, цитаты и т.п.), и объединяются в единый текст. После извлечения текст дополнительно очищается (удаляются сноски вида [1], [2–5]; восстанавливаются пробелы после знаков препинания; исправляются слипшиеся слова между кириллицей и латиницей, а также внутри CamelCase)
\end{itemize}

\textbf{- Обработка ошибок}

Перед сохранением документа выполняется валидация HTML
страницы на предмет ошибок или отсутствия полезного содержимого. Эта 
логика реализована в функции is\_invalid\_page. Функция проверяет:
\begin{itemize}
    \item наличие фразы «404» или «страница не найдена» в тегах <title> и <h1>
    \item общий объём текста на странице (если менее 10 слов — страница 
    считается пустой или ошибочной) 
\end{itemize}
Если is\_invalid\_page возвращает True, документ не сохраняется, и 
робот переходит к следующему URL. 

Также при HTTP-запросе обрабатывается следующее: 
\begin{itemize}
    \item статус 404 — явно пропускается с сообщением
    \item любой статус, отличный от 200, также приводит к пропуску
    \item исключения при сетевом запросе (таймаут, DNS-ошибка и т.п.) 
    перехватываются и логируются
\end{itemize}
Благодаря данным мерам обеспечивается высокое качество корпуса: 
в базу попадают только валидные, полные и содержательные статьи. 

\subsection*{Выводы}

В результате выполнения данной лабораторной был реализован 
поисковой робот, собирающий корпус документов по выбранной теме. 
Были учтены структуры обоих источник при парсинге, была корректно 
собрана мета-информация и выделен чистых текст статей, которые 
сохраняются в базу данных. Также поисковой робот способен обновлять 
корпус документов. В случае остановок состояние сохраняется, все ошибки 
успешно обрабатываются. Собираются только валидные статьи.

\pagebreak

\section{Лабораторная работа \textnumero 3 \enquote{Токенизация}}

Нужно реализовать процесс разбиения текстов 
документов на токены, который потом будет использоваться при 
индексации. Для этого потребуется выработать правила, по которым текст 
делится на токены. Необходимо описать их в отчёте, указать достоинства и недостатки выбранного метода. Привести примеры токенов, которые были 
выделены неудачно, объяснить, как можно было бы поправить правила, 
чтобы исправить найденные проблемы. 

В результатах выполнения работы нужно указать следующие 
статистические данные: 
\begin{itemize}
    \item Количество токенов
    \item Среднюю длину токена
\end{itemize}
Кроме того, нужно привести время выполнения программы, указать 
зависимость времени от объёма входных данных. Указать скорость 
токенизации в расчёте на килобайт входного текста. Является ли эта 
скорость оптимальной? Как её можно ускорить?

\subsection*{Описание}

Токенизация — это фундаментальный процесс в компьютерной 
лингвистике и обработке естественного языка, заключающийся в разбиении 
непрерывного текста на минимальные смысловые единицы — токены. В 
широком смысле токеном может считаться слово, число, пунктуационный 
знак или даже отдельная морфема в зависимости от поставленной задачи и 
выбранного уровня гранулярности. 

\textbf{- Выбранные правила токенизации} 
\begin{itemize}
    \item Поддержка кириллицы: токены могут содержать русские буквы 
    \item Допустимые символы: буквы (латинские и кириллические), цифры 
    (только в составе специальных числовых токенов), дефис (только 
    внутри слова, не в начале/конце)
    \item Фильтрация по длине: максимальная длина токена — 20 символов, 
    минимальное требование для слов с дефисом: каждая часть должна 
    содержать >= 2 кириллических символа
    \item Отсев некорректных последовательностей: исключаются токены с 
    тройным повторением одной буквы (например, "аааа"), запрещаются 
    токены, состоящие только из дефисов или цифр (кроме отдельных 
    чисел) 
    \item Нормализация регистра — все токены приводятся к нижнему регистру для унификации 
    \item Список разрешенных аббревиатур — специальный файл abbrevs.txt 
    содержит исключения, которые принимаются без дополнительной 
    проверки 
    \item Числовые токены — допускаются только: отдельные числа (например, "2023"), максимальная длина — 4 цифры, запрещены числа с ведущими нулями (кроме "0") 
\end{itemize}

\subsection*{Ход работы}

\textbf{- Принцип работы} 
\begin{itemize}
\item Предобработка текста (текст ограничивается первыми 100 000 символов для оптимизации) 
\item Построение токенов (текст сканируется байт за байтом, 
последовательности допустимых символов накапливаются в буфер) 
\item Фильтрация и нормализация (удаление дефисов в начале/конце буфера, проверка по списку разрешенных аббревиатур, приведение к нижнему регистру)
\item Валидация токена (для русского слова: хотя бы одна кириллица, длина <= 20 символов, не содержит тройных повторов букв, при наличии 
дефиса — каждая часть >= 2 кириллических символа) 
\item Сохранение результатов (каждому документу присваивается ID, токены сохраняются в файлы tokens/{id}.tokens) 
\end{itemize}

\textbf{- Результаты} 

\begin{center}
\begin{tabular}{ |c|c| }
\hline
%Всего токенов & 20447144 \\
 Всего токенов & 37487814 \\  
 \hline
 %Средняя длина токена & 6.24 символов \
 Средняя длина токена & 6.16 символов \\ 
 \hline
 %Время выполнения & 32.63 сек \\
 Время выполнения & 55.23 сек \\
 \hline
 %Скорость токенизации & 8700.43 КБ/сек \\
 Скорость токенизации & 9299.27 КБ/сек \\
 \hline
\end{tabular}
\end{center} 

\subsection*{Выводы}
Статистика в результате показывает, что алгоритм обеспечивает 
достаточно плотную токенизацию (средняя длина примерно 6 символов), что соответствует типичной длине русских слов. Высокая скорость обработки (55.23 секунды на весь корпус) подтверждает эффективность реализации для крупномасштабной обработки текстов.

\pagebreak

\section{Лабораторная работа \textnumero 4 \enquote{Стемминг}}

Добавить в созданную поисковую систему лемматизацию. 
В простейшем случае, это просто поиск без учёта словоформ. В более 
сложном случае, можно давать бонус большего размера за точное 
совпадение слов. 

Стемминг можно добавлять на этапе индексации, можно на этапе 
выполнения поискового запроса. В отчёте должна быть включена оценка 
качества поиска, после внедрения стемминга. 


\subsection*{Описание}

Стемминг (stemming — выделение основы) — это процесс 
лингвистической нормализации текста, при котором слова приводятся к 
своей корневой форме (основе) путем отсечения словоизменительных и 
словообразовательных аффиксов. 

Основные задачи стемминга: 
\begin{itemize}
\item Уменьшение размерности пространства признаков в задачах NLP; 
\item Группировка семантически родственных слов (например: "читает", 
"читающий", "читатель" → "чита") 
\item Повышение релевантности поиска за счет учета различных форм слова 
\item Упрощение анализа текстовой статистики 
\end{itemize}
В отличие от лемматизации, стемминг не обязательно возвращает 
словарную форму слова, а выделяет общую основу для группы родственных 
слов. 


\subsection*{Ход работы}

\textbf{- Принцип работы} 
\begin{itemize}
\item Предобработка токенов (алгоритм получает на вход уже готовые токены, которые были ранее извлечены из текстов и сохранены в файлах 
формата .tokens, каждый токен обрабатывается независимо) 
\item Фильтрация стоп-слов (перед стеммингом каждый токен проверяется по загруженному списку стоп-слов (из файла stop\_words\_ru.txt), если токен совпадает со стоп-словом, он полностью исключается из дальнейшей 
обработки, что сокращает шум и уменьшает объем данных)
\item Построение стем (для каждого токена начала проверяются исключения 
(короткие слова и неизменяемые глаголы), затем специальные правила 
замены (например, "ирование" → "ир"), далее обработка возвратных 
глаголов, и только в последнюю очередь — стандартные суффиксы из 
предварительно отсортированного по длине списка) 
\item После каждой трансформации проверяется, что полученная основа 
содержит не менее 2 Unicode-символов. Более короткие результаты 
отбрасываются, что предотвращает появление неинформативных 
фрагментов 
\end{itemize}

\textbf{- Результаты} 

\begin{center}
\begin{tabular}{ |c|c| }
\hline
%Всего стемм & 13173177 \\
 Всего стемм & 23910847 \\  
 \hline
 %Отброшено & 7273967 \\
 Отброшено & 13576967 \\ 
 \hline
 %Время работы & 49.10 сек \\
 Время работы & 75.85 сек \\
 \hline
 %Средняя скорость & 5.32 МБ/сек \\
 Средняя скорость & 6.25 МБ/сек \\
 \hline
\end{tabular}
\end{center} 

\subsection*{Выводы}
Статистика обработки показывает высокую эффективность 
алгоритма: за 75.85 секунд были обработаны все файлы. Отброшено 13.58 миллионов токенов (36.2\% от общего количества), что подтверждает качественную фильтрацию стоп-слов и неинформативных коротких основ. Получено 23.91 миллионов стемов со средней скоростью 6.25 МБ/сек, что указывает на хорошо сбалансированный алгоритм. Подход стеммера (сохранение коротких слов без изменений, минимальная длина основы 2 символа) обеспечивает сохранение семантической ценности, избегая излишней обрезки. 

\pagebreak

\section{Лабораторная работа \textnumero 5 \enquote{Закон Ципфа}}

Для своего корпуса необходимо построить график распределения терминов по частотностям в логарифмической шкале, наложить на этот график закон Ципфа. Объяснить причины расхождения. 

В качестве дополнительного задания, можно (но необязательно) 
подобрать константы для закона Мандельброта, наложить полученный 
график на график распределения терминов по частотностям. Привести 
выбранные константы. 

\subsection*{Описание}

Закон Ципфа (Zipf's Law) — это эмпирическая закономерность, 
описывающая распределение частот слов в естественных языках. 

Если упорядочить слова текста по убыванию частоты, то частота $r$-го 
по популярности слова примерно обратно пропорциональна его рангу $r$:
$f(r) = \frac{C}{r}$,

где:

$f(r)$ — частота слова с рангом $r$, 

$r$ — ранг слова (1 для самого частотного), 

$C$ — константа, примерно равная частоте самого частого слова.

Закон Ципфа помогает проверить, обладает ли исследуемый набор 
текстов фундаментальными свойствами естественного языка. Благодаря 
ему можно определить, что очень небольшое количество слов используется 
чрезвычайно часто, выполняя грамматическую и связующую функцию, в то 
время как подавляющее большинство слов встречается редко, обеспечивая 
смысловое разнообразие и конкретику.

\subsection*{Ход работы}

Для анализа статистических свойств корпуса научно-популярных 
статей был построен график распределения частот токенов в 
логарифмических координатах (log-log шкала) и наложена теоретическая 
кривая, соответствующая закону Ципфа. 

На графике наблюдается характерное распределение, где:
\begin{itemize}
    \item фактические частоты токенов (синие точки) образуют кривую, близкую к прямой линии в логарифмических координатах
    \item теоретический закон Ципфа (красная линия) представляет собой 
идеальную прямую с наклоном -1
    \item синяя кривая располагается несколько выше теоретической прямой
\end{itemize}

\begin{figure}[H]
    \centering
   \includegraphics[width=1\linewidth]{zipf.png}
    \caption{Закон Ципфа}
    \label{fig:placeholder}
\end{figure}

%\begin{figure}[H]
%    \centering
%    \includegraphics[width=0.9\linewidth]{zipf_stem.png}
%    \caption{Закон Ципфа по стеммам}
%    \label{fig:placeholder}
%\end{figure}

\subsection*{Выводы}

Данные в целом следуют закону Ципфа.  
\begin{itemize}
    \item в начале (низкий ранг, высокая частота): самые частотные слова (ранги 1-5) имеют частоту, которая выше, чем предсказывает идеальный закон Ципфа. Это типичное явление, так как самые частотные функциональные слова (например, "и", "в", "на", "не") встречаются чаще
    \item в конце (высокий ранг, низкая частота): При очень высоких рангах (редкие слова) фактические частоты также могут отклоняться от теоретической линии, но они потихоньку приближаются к ней
\end{itemize}
Однако наблюдаются расхождения, которые могут быть из-за того, что научно-популярные тексты более стандартизированы, чем естественный язык. Также такое возможно из-за часто повторяющихся конструкций, клише, стандартных фраз, характерных для научно-популярного стиля.

\pagebreak

\section{Лабораторная работа \textnumero 6 \enquote{Булев индекс}}

Необходимо построить булев индекс для дальнейшего его использования при поиске.

\subsection*{Описание}
Булев индекс (или инвертированный индекс) — это основная структура данных, используемая в системах информационного поиска. Он представляет собой отображение термин → список идентификаторов документов, в которых этот термин встречается (так называемый posting list). Такой индекс позволяет эффективно выполнять булевы запросы, такие как:
\begin{itemize}
    \item термин1 AND термин2 — документы, содержащие оба термина
    \item термин1 OR термин2 — документы, содержащие хотя бы один из терминов
    \item NOT термин — документы, не содержащие термин.
\end{itemize}
Индекс строится следующим образом:
\begin{itemize}
    \item Производится токенизация текста (разбиение на слова или термины)
    \item Применяются фильтрация (удаление стоп-слов) и, при необходимости, нормализация (лемматизация, стемминг)
    \item Для каждого документа фиксируется, какие термины в нём присутствуют (часто без учёта частоты — только факт наличия)
    \item Собираются все термины в глобальный словарь, и для каждого термина формируется список ID документов
    \item Словарь и posting-листы сохраняются на диск, обычно в отсортированном виде — это ускоряет последующие поисковые операции
\end{itemize}

\subsection*{Ход работы}

\textbf{- Построение} 

Индекс строится на основе предварительно обработанных файлов, хранящихся в директории ../../preprocess/tokenization/tokens. 
Каждый файл соответствует одному документу и содержит по одной строке на термин.

Этапы:
\begin{itemize}
    \item Чтение данных: каждый .tokens-файл читается как список терминов. Файлы обрабатываются в цикле по директории, а ID документа извлекается из имени файла (например, 12345.stems → doc\_id = 12345)
    \item Удаление дубликатов на уровне документа: функция deduplicate\_terms() сортирует термины вставкой и удаляет повторяющиеся — это гарантирует, что один документ не будет многократно добавлен в posting list одного термина
    \item Построение индекса: используется собственная хэш-таблица (TermIndex) с открытой адресацией для отслеживания, какие термины уже встречались. При первом вхождении термина он добавляется в глобальный вектор dictionary, а в inverted\_lists создаётся новый posting list с текущим doc\_id. При повторном вхождении термина (в другом документе) doc\_id добавляется в уже существующий список
    \item Сортировка и дедупликация posting lists: после обработки всех документов термины сортируются лексикографически вместе с их списками (sort\_lexicographically). Каждый posting list сортируется (вставками) и очищается от возможных дубликатов — хотя они теоретически не должны появляться, это страховка на случай ошибок на более ранних этапах
    \item Сохранение: индекс записывается в текстовый файл в формате термин:doc\_id1,...
\end{itemize}

\textbf{- Результаты} 

Построение индекса по токенам:

\begin{center}
\begin{tabular}{ |c|c| }
\hline
 %Терминов & 522389 \\ 
 Терминов & 679469 \\
 \hline
 %Время & 3851.96 сек \\ 
 Время & 6962.59 сек \\
 \hline
 %Объем & 86.95 МБ \\ 
 Объем & 160.22 МБ \\
 \hline
\end{tabular}
\end{center} 

Дополнительно еще есть построение индекса по стеммам:

\begin{center}
\begin{tabular}{ |c|c| }
\hline
 %Терминов & 344344 \\ 
 Терминов & 452372 \\
 \hline
 Время & 2773.18 сек \\
 %Время & 1567.28 сек \\ 
 \hline
 Объем & 114.61 МБ \\
 %Объем & 62.57 МБ \\ 
 \hline
\end{tabular}
\end{center} 

\subsection*{Выводы}
Стемминг сократил словарь на ~34\% (с 522 тысяч до 452 тысяч терминов), что ожидаемо: формы одного слова (например, рисует, рисовал, рисование) сводятся к одному стему. Время построения сократилось более чем в 2.4 раза. Это объясняется не только меньшим числом терминов, но и более короткими строками (стемы короче полных слов), что ускоряет хэширование и сравнение.

\pagebreak

\section{Лабораторная работа \textnumero 7 \enquote{Булев поиск}}

Необходимо написать булев поиск как утилиту командной строки, получающую запросы со стандартного файла ввода и выдающую результат в стандартный файл вывода.

\subsection*{Описание}
Булев поиск — это классический метод информационного поиска, основанный на логических операциях И (AND), ИЛИ (OR) и НЕ (NOT).
\begin{itemize}
    \item Запрос строится как логическое выражение над терминами.
    \item Результат — множество документов, в которых выполняется это выражение.
    \item Порядок релевантности отсутствует: все найденные документы считаются равнозначными.
\end{itemize}

\subsection*{Ход работы}

\textbf{- Как работает}

\begin{itemize}
    \item Парсинг запроса: пользователь вводит запрос в виде строки (например, нейрон, нейрон and память, нейрон and not рак)
    \item Индекс: загружается из файла boolean\_index.txt, где каждая строка имеет вид: термин:doc1,doc2,doc3,...
    \item Хранение: индекс представлен в виде std::vector<IndexedTerm>, где термины отсортированы лексикографически
    \item Поиск термина: выполняется бинарным поиском по отсортированному вектору
    \item Обработка запроса: 
    \begin{itemize}
        \item Для одного термина — возвращается его posting list
        \item Для составных запросов (A and B, A or B, A and not B) — выполняются соответствующие операции над списками (AND — compute\_intersection, OR — compute\_union, NOT — compute\_difference)
    \end{itemize}
\end{itemize}

\textbf{- Результаты поиска}

\begin{minted}{text}
нейрон
Найдено: 198
[ID: 4789] — https://www.techinsider.ru/science/6788-po-obrazu-i-podobiyu-iskusstvennyy-razum/
[ID: 5277] — https://www.techinsider.ru/science/7286-mikro-mysli-sila-neyrona/
[ID: 5326] — https://www.techinsider.ru/science/7335-siloy-mysli-obezyana-upravlyaet-robotom/
[ID: 6503] — https://www.techinsider.ru/technologies/8622-pervoe-slovo-skanirovanie-mozga/
[ID: 6958] — https://www.techinsider.ru/technologies/9078-drevneyshee-iz-chuvstv-zapakhi/
[ID: 8107] — https://www.techinsider.ru/science/10229-partiya-mozga-neyronnoe-soobshchestvo/
[ID: 8329] — https://www.techinsider.ru/technologies/10451-zapakh-sinego-chem-pakhnet-tsvet/
[ID: 8482] — https://www.techinsider.ru/science/10606-mukhi-v-virtualnoy-realnosti-tonkosti-raspoznavaniya/
[ID: 9062] — https://www.techinsider.ru/science/11188-bezdele-utomlyaet-energiya-v-nikuda/
[ID: 9159] — https://www.techinsider.ru/science/11285-10-tekhnicheskikh-kontseptsiy-o-kotorykh-stoit-znat-v-2011-godu/
(Время поиска: 0.03 мс)
\end{minted}

\begin{minted}{text}
прогресс and телефон
Найдено: 23
[ID: 4766] — https://www.techinsider.ru/technologies/6765-arsenal-supershpiona-dzheyms-bond-i-ego-shtuchki/
[ID: 5005] — https://www.techinsider.ru/technologies/7009-kinoteatr-na-nosu-blizhniy-obzor/
[ID: 8970] — https://www.techinsider.ru/technologies/11095-10-proryvov-desyatiletiya-innovatsii-2020/
[ID: 9149] — https://www.techinsider.ru/technologies/11275-novosti-kosmonavtiki-21-27-yanvarya-2011-g/
[ID: 9450] — https://www.techinsider.ru/technologies/11577-novosti-kosmonavtiki-22-28-aprelya-2011-g/
[ID: 10414] — https://www.techinsider.ru/technologies/12541-neytrinnaya-perepiska-slovo-skvoz-kamen/
[ID: 12577] — https://biomolecula.ru/articles/kompiuternye-igry-v-molekuliarnuiu-biofiziku-biologicheskikh-membran
[ID: 13250] — https://biomolecula.ru/articles/tema-pesni-dlia-tsoia-karl-landshteiner
[ID: 26819] — https://www.techinsider.ru/made-in-russia/275122-rossiyskie-uchenye-zaryadili-mobilnyy-s-1-5-kilometrov/
[ID: 30110] — https://www.techinsider.ru/weapon/352762-gibridnye-smart-chasy-novoe-napravlenie-chasovogo-dela-na-primere-casio-edifice-eqb-501xdb/
(Время поиска: 0.42 мс)
\end{minted}

\begin{minted}{text}
технология or айфон
Найдено: 4300
[ID: 3311] — https://www.techinsider.ru/science/5278-nebesnyy-glaz-kosmicheskomu-teleskopu-khabbl-ispolnilos-pyatnadtsat-let/
[ID: 3325] — https://www.techinsider.ru/adrenalin/5292-chernyy-belyy-i-nemnogo-zolota-i-kruglee-vsekh-predshestvennikov/
[ID: 3332] — https://www.techinsider.ru/technologies/5299-futbol-kak-dvigatel-progressa-igra-vysokogo-razresheniya/
[ID: 3335] — https://www.techinsider.ru/technologies/5302-ledenyashchie-podrobnosti-poklonenie-snezhnoy-koroleve/
[ID: 3336] — https://www.techinsider.ru/technologies/5303-krasota-na-odnu-noch-odnorazovaya-posuda-klassa-de-luxe/
[ID: 3341] — https://www.techinsider.ru/science/5308-nasledie-nobelya-2005-i-novye-nasledniki/
[ID: 3342] — https://www.techinsider.ru/technologies/5309-e-t-razbushevalsya-iskusstvo-modelirovat-koshmary/
[ID: 3361] — https://www.techinsider.ru/design/5329-ego-zovut-kong-king-kong-kak-obezyana-proizoshla-ot-cheloveka/
[ID: 3362] — https://www.techinsider.ru/gadgets/5330-ubiytsy-bolshikh-televizorov-probuem-tovarishchey-na-tsvet/
[ID: 3381] — https://www.techinsider.ru/history/5350-koronovannaya-gazirovka-istoriya-pivnoy-probki/
(Время поиска: 0.07 мс)
\end{minted}

\begin{minted}{text}
нейролептик and not приступ
Найдено: 2
[ID: 14207] — https://biomolecula.ru/articles/kratkaia-istoriia-antidepressantov
[ID: 32466] — https://www.techinsider.ru/science/news-382882-ustanovleno-proishozhdenie-shizofrenii/
(Время поиска: 0.04 мс)
\end{minted}

\subsection*{Выводы}
В итоге, был реализован быстрый булев поиск для точных логических запросов.

\pagebreak

\section{Лабораторная работа \textnumero 8 \enquote{Сжатие}}

Необходимо применить к построенному булеву индексу сжатие и провести замеры.

\subsection*{Описание}
Сжатие инвертированного индекса применяется для снижения объёма занимаемой памяти и ускорения операций поиска за счёт уменьшения количества чтений с диска. Инвертированный индекс содержит для каждого термина список идентификаторов документов (posting list), в которых он встречается. Эти списки часто обладают высокой избыточностью:

Идентификаторы документов, как правило, упорядочены по возрастанию.
Разности (дельты) между соседними значениями значительно меньше самих значений.

Основные подходы к сжатию:
\begin{itemize}
    \item Delta-кодирование преобразует список документов в последовательность разностей между соседними ID, что уменьшает разрядность чисел
    \item Переменное побайтовое (Variable Byte, VB) кодирование представляет целые числа в компактной форме, используя от 1 до 5 байт в зависимости от величины значения
    \item Front Coding сжимает термины, сохраняя только суффикс текущего термина относительно предыдущего, если у них есть общий префикс
\end{itemize}

\subsection*{Ход работы}

\textbf{- Примененные методы}
\begin{itemize}
    \item Delta-кодирование: перед сжатием каждый posting list преобразуется в последовательность разностей $\Delta_0=d_0, \Delta_i=d_i-d_{i-1} (i>0)$, где $d_i$ —  ID документа в исходном списке. Это приводит к значительному уменьшению значений, особенно в больших корпусах с плотной нумерацией документов
    \item VB-кодирование: каждое число из последовательности дельт кодируется методом Variable Byte:
    \begin{itemize}
        \item число разбивается на 7-битные блоки
        \item каждый блок записывается в отдельный байт, где старший (8-й) бит указывает, есть ли продолжение (1 — продолжение, 0 — конец числа)
        \item это позволяет компактно кодировать как малые, так и большие значения без фиксированного размера
    \end{itemize}
    \item Front Coding терминов: для сжатия строк терминов используется front coding:
    \begin{itemize}
        \item для текущего термина определяется длина общего префикса с предыдущим термином
        \item записывается длина префикса и суффикс (оставшаяся часть)
        \item если длина префикса или суффикса превышает 255 байт, используется «fallback» — термин сохраняется целиком с маркером 0xFF
    \end{itemize}
\end{itemize}
Все данные записываются в бинарный файл с чёткой структурой записи для каждого термина: [suffix\_len:1][prefix\_len:1][suffix][data\_len:4][VB-кодированные дельты].

\textbf{- Результаты} 

Сжатие индекса:

\begin{center}
\begin{tabular}{ |c|c| }
\hline
 %Исходный размер & 86.95 МБ \\ 
 Исходный размер & 160.22 МБ \\
 \hline
 %Сжатый размер & 23.1 МБ \\ 
 Сжатый размер & 40.17 МБ \\
 \hline
\end{tabular}
\end{center} 

Дополнительно еще есть сжатие индекса, построенного по стеммам:

\begin{center}
\begin{tabular}{ |c|c| }
\hline
 %Исходный размер & 62.57 МБ \\ 
 Исходный размер & 114.61 МБ \\
 \hline
 %Сжатый размер & 16.5 МБ \\ 
 Сжатый размер & 28.66 МБ \\
 \hline
\end{tabular}
\end{center} 

\subsection*{Выводы}
В результате объём данных сократился почти в 4 раза, что значительно снижает требования к памяти и дисковому пространству.

\pagebreak

\section{Лабораторная работа \textnumero 9 \enquote{TF-IDF}}

Необходимо построить частотный индекс для поиска с TF-IDF.

\subsection*{Описание}
TF-IDF (Term Frequency – Inverse Document Frequency) — это статистическая мера, используемая для оценки важности термина в документе относительно всего корпуса документов. Она широко применяется в задачах информационного поиска, ранжирования релевантных документов и векторного представления текста.

Формула TF-IDF для термина $t$ в документе $d$: $TF$-$IDF(t,d)=TF(t,d)*IDF(t)$, где:
\begin{itemize}
    \item $TF$ (Term Frequency) — частота встречаемости термина $t$ в документе $d$ (я использовала количество вхождений)
    \item $IDF$ (Inverse Document Frequency) — мера «редкости» термина в корпусе: $IDF(t) = log(\frac{N}{df_t})$, где
    \begin{itemize}
        \item $N$ — общее число документов в корпусе
        \item $df_t$ — количество документов, содержащих термин $t$
    \end{itemize}
\end{itemize}
Чем чаще термин встречается в конкретном документе и чем реже — в других документах, тем выше его $TF$-$IDF$. Это позволяет:
\begin{itemize}
    \item уменьшить вес общих (стоп-)слов (например, «и», «в»), которые имеют высокий $TF$, но также высокий $IDF$
    \item увеличить вес терминов, характерных именно для данного документа (например, «барокко» в статье об искусстве)
\end{itemize}

Ранжированный (TF-IDF) поиск поиск оценивает степень релевантности каждого документа по отношению к запросу.

Используется мера TF-IDF:
\begin{itemize}
    \item TF (Term Frequency) — частота термина в документе
    \item IDF (Inverse Document Frequency) — редкость термина в корпусе
\end{itemize}
Итоговый score документа — сумма TF-IDF значений всех терминов запроса, встречающихся в нём.

Результат сортируется по убыванию score, что обеспечивает ранжирование по релевантности.

\subsection*{Ход работы}

\textbf{- Построение} 
Выполнялось построение двух файлов, необходимых для последующего вычисления или использования TF-IDF:

\begin{itemize}
    \item Файл с частотным индексом ($TF$-index): Для каждого термина сохраняется список пар (doc\_id, term\_frequency). Результат вв иде: термин:doc1:freq1,doc2:freq2,...
    \item Файл $IDF$: для каждого термина вычисляется значение $IDF$ по формуле. Результат в виде: термин значение\_idf
\end{itemize}

В результает было получено два файла:
\begin{itemize}
    \item tf\_index.txt — частотный индекс 
    \item idf.txt — значения $IDF$ для каждого термина
\end{itemize}

\textbf{- Результаты построения} 

Программа была запущена отдельно для токенов и стемм:

\begin{center}
\begin{tabular}{ |c|c| }
\hline
 Токены & 679469 уникальных терминов \\ 
 \hline
 Стеммы & 452372 уникальных терминов \\ 
 \hline
\end{tabular}
\end{center} 

\textbf{- Поиск}

Исходные данные:
\begin{itemize}
    \item idf.txt — файл с IDF-весами: термин значение\_idf
    \item tf\_index.txt — файл с частотами: термин:doc1:freq1,doc2:freq2,...
    \item document\_lengths — вычисляются из tf\_index.txt как сумма всех частот терминов в документе
\end{itemize}

Вычисление score:
\begin{itemize}
    \item для каждого термина запроса берётся IDF.
    \item для каждого документа, где термин встречается, вычисляется нормализованный $TF$: $TF$ = частота термина/общее число терминов в документе
    \item score документа — сумма $TF$ * $IDF$ по всем терминам запроса.
\end{itemize}

Ранжирование: документы сортируются по убыванию score.

Вывод: до 10 документов с указанием score и URL.

\textbf{- Результаты поиска}

\begin{minted}{text}
технический
Найдено: 549
[ID: 53942] (score: 0.0958) — https://www.techinsider.ru/weapon/news-608353-su-57-podelitsya-s-tu-160m2-i-tu-22m3m-sistemoy-svyazi/
[ID: 96081] (score: 0.0893) — https://www.techinsider.ru/popmem/1630577-kak-po-cvetu-truby-opredelit-chto-u-nee-vnutri/
[ID: 16643] (score: 0.0857) — https://www.techinsider.ru/technologies/15498-samsung-electronics-voshla-v-konsortsium-openpower-foundation/        
[ID: 43562] (score: 0.0732) — https://www.techinsider.ru/editorial/502222-independent-media-primet-uchastie-v-big-dataandai-conference/
[ID: 41963] (score: 0.0693) — https://www.techinsider.ru/technologies/news-485012-40-indiycev-zastryali-na-karuseli/
[ID: 61103] (score: 0.0693) — https://www.techinsider.ru/technologies/news-679903-michelin-prevratit-kukuruzu-i-ris-v-avtomobilnye-pokryshki/     
[ID: 33663] (score: 0.0667) — https://www.techinsider.ru/adrenalin/395332-kak-rabotaet-aerodinamika-v-formule-1/
[ID: 32941] (score: 0.0667) — https://www.techinsider.ru/technologies/news-387802-polzovateli-ps4-i-xbox-one-vpervye-smogli-sygrat-mezhdu-soboy/  
[ID: 40084] (score: 0.0659) — https://www.techinsider.ru/vehicles/464452-unikalnuyu-mazda-rx-7-nashli-spustya-35-let/
[ID: 47103] (score: 0.0606) — https://www.techinsider.ru/weapon/news-538814-sovetskiy-ulyanovsk-pomozhet-v-sozdanii-novogo-avianosca-dlya-vmf-rossii/
(Время поиска: 0.25 мс)
\end{minted}

\begin{minted}{text}
нейрон
Найдено: 198
[ID: 81255] (score: 0.0967) — https://www.techinsider.ru/technologies/news-1565145-sozdan-iskusstvennyy-neyron-kotoryy-mozhno-soedinit-s-zhivoy-nervnoy-kletkoy/
[ID: 81259] (score: 0.0762) — https://www.techinsider.ru/popmem/1565173-perelivanie-sinteticheskoy-krovi-i-iskusstvennyy-neyron-kotorye-soedinyaetsya-s-zhivym-analogom-glavnye-novosti-segodnya/
[ID: 12010] (score: 0.0760) — https://www.techinsider.ru/science/14145-vpolsily-mozg-i-virtualnaya-realnost/
[ID: 54346] (score: 0.0709) — https://www.techinsider.ru/technologies/news-612433-supersposobnost-saranchi-adaptirovali-dlya-bespilotnikov/       
[ID: 107969] (score: 0.0690) — https://www.techinsider.ru/news/news-1682853-genetiki-prevratili-kletki-koji-v-neirony-dlya-kletochnoi-terapii-mozga/
[ID: 5277] (score: 0.0645) — https://www.techinsider.ru/science/7286-mikro-mysli-sila-neyrona/
[ID: 12778] (score: 0.0485) — https://biomolecula.ru/articles/robert-sapolski-vsio-resheno-zhizn-bez-svobody-voli-retsenziia
[ID: 92173] (score: 0.0464) — https://www.techinsider.ru/technologies/news-1613507-razrabotan-integrirovannyy-neyron-kotoryy-mozhet-odnovremenno-videt-i-osyazat/
[ID: 61105] (score: 0.0452) — https://www.techinsider.ru/editorial/679923-kakoy-obem-pamyati-u-nashego-mozga-i-mozhet-li-ona-zakonchitsya/        
[ID: 82972] (score: 0.0417) — https://www.techinsider.ru/editorial/1573153-mozhno-li-uvelichit-vyrabotku-serotonina-samostoyatelno/
(Время поиска: 0.13 мс)
\end{minted}

\subsection*{Выводы}
В результате был построен частотный индекс и вычислены IDF значения для каждого термина, что позволило реализовать ранжированный поиск: документы теперь упорядочиваются по релевантности запросу на основе количественной оценки важности термина в каждом конкретном документе относительно всего корпуса.

\pagebreak

% \section{Лабораторная работа \textnumero 10 \enquote{Тесты}}

% Написать тесты.

% \subsection*{Описание}

% \subsection*{Ход работы}

% \subsection*{Выводы}

\end{document}